% =============================================================================
%  Zdrojové texty – Záhrada stratených ciest
%  Autor: Mykyta Olym, Počítačová grafika 2026/2027
% =============================================================================
\documentclass[11pt, a4paper]{article}
\newcommand{\druhprirucky}{Zdrojové texty}
% =============================================================================
%  Spoločné nastavenie všetkých troch dokumentov (preklad: latexmk -xelatex)
%  Záhrada stratených ciest – izometrické bludisko v knihe, Mykyta Olym, PG 2026/2027
% =============================================================================

% ---------------------------------------------------------------- údaje projektu
\newcommand{\nazovprojektu}{Záhrada stratených ciest}
\newcommand{\podnazovprojektu}{izometrické bludisko v knihe}
\newcommand{\predmet}{Počítačová grafika}
\newcommand{\skolskyrok}{2026/2027}
\newcommand{\ucitel}{Course Instructor}
\newcommand{\autor}{Mykyta Olym}
\newcommand{\verzia}{2.0}

% ---------------------------------------------------------------- jazyk a písmo
\usepackage{fontspec}
\usepackage{polyglossia}
\setdefaultlanguage{slovak}
\setmainfont{TeX Gyre Pagella}
\setsansfont{TeX Gyre Heros}
% neproporcionálne písmo zo šablóny KPI (priečinok fonts/)
\setmonofont{SourceCodePro}[
  Path = fonts/, Extension = .otf, Scale = 0.85,
  UprightFont = *-Regular, BoldFont = *-Bold,
  ItalicFont = *-Italic, BoldItalicFont = *-BoldItalic]

% ---------------------------------------------------------------- strana
\usepackage[a4paper, top=25mm, bottom=30mm, left=28mm, right=25mm]{geometry}
\usepackage{fancyhdr}
\pagestyle{fancy}
\fancyhf{}
\fancyhead[L]{\small\nazovprojektu}
\fancyhead[R]{\small\druhprirucky}
\fancyfoot[R]{\thepage}
\renewcommand{\headrulewidth}{0.4pt}
\setlength{\headheight}{14pt}
\setlength{\parindent}{0pt}
\setlength{\parskip}{0.55em}
\setlength{\emergencystretch}{3em}   % dlhé názvy radšej zalomiť než vytlačiť za okraj
% názvy súborov sa smú zalomiť za lomkou, bodkou alebo podčiarknutím
\usepackage[hyphens]{url}
\newcommand{\subor}[1]{\nolinkurl{#1}}

% ---------------------------------------------------------------- obsah, obrázky, tabuľky
\usepackage{graphicx}
\graphicspath{{obrazky/}}
\usepackage{booktabs, tabularx, array, longtable}
\usepackage{float}
\usepackage[font=small, labelfont=bf]{caption}
\usepackage{enumitem}
\setlist{itemsep=0.15em, topsep=0.3em}
\usepackage{xcolor}
\usepackage{csquotes}
\usepackage[style=iso-numeric, backend=biber, sorting=none]{biblatex}
\addbibresource{literatura.bib}
\usepackage[hidelinks]{hyperref}

% ---------------------------------------------------------------- titulná strana podľa pokynov
\newcommand{\titulnastrana}{%
  \begin{titlepage}
    \centering
    {\large Katedra počítačov a informatiky\\
      Fakulta elektrotechniky a informatiky\\
      Technická univerzita v Košiciach\par}
    \vspace*{\fill}
    {\Huge\bfseries \nazovprojektu\par}
    \vspace{0.4em}
    {\large \podnazovprojektu\par}
    \vspace{1.2em}
    {\Large \predmet\par}
    \vspace{3em}
    {\LARGE \druhprirucky\par}
    \vspace*{\fill}
    \begin{minipage}[t]{0.62\textwidth}
      \raggedright
      \begin{tabular}{@{}ll@{}}
        Šk. rok:       & \skolskyrok \\
        Meno učiteľa:  & \ucitel
      \end{tabular}
    \end{minipage}%
    \begin{minipage}[t]{0.38\textwidth}
      \raggedleft
      \begin{tabular}{@{}ll@{}}
        Autori: & \autor
      \end{tabular}
    \end{minipage}
  \end{titlepage}
  \setcounter{page}{1}%
}

% Obsah je podľa pokynov kapitola 1, preto skutočné kapitoly začínajú číslom 2.
\newcommand{\obsah}{%
  \renewcommand{\contentsname}{1\quad Obsah}%
  \tableofcontents
  \setcounter{section}{1}%
  \clearpage
}


\newcommand{\hra}{../Olym_PG1_20262027}
\newcommand{\nastroje}{../nastroje}

% Výpisy cez fvextra: na rozdiel od balíka listings správne vysádza všetky znaky Unicode
% (č, ž, ľ, š, šípky) a dlhé riadky zalomí so značkou pokračovania.
\usepackage{fvextra}
\fvset{numbers=left, numbersep=6pt, fontsize=\footnotesize, frame=lines, framerule=0.4pt,
  breaklines=true, breakanywhere=true, tabsize=4, xleftmargin=1.8em}
\renewcommand{\theFancyVerbLine}{\tiny\color{black!50}\arabic{FancyVerbLine}}

% Výpis jedného súboru ako samostatná podkapitola: koreň, cesta, popis
\newcommand{\vypis}[3]{%
  \subsection{\texorpdfstring{\subor{#2}}{#2}}
  #3
  \VerbatimInput{#1/#2}%
}

\begin{document}
\titulnastrana
\obsah

% =============================================================================
\section{Úvod}

Dokument obsahuje úplné zdrojové texty programu Záhrada stratených ciest verzie \verzia{} a nástrojov,
ktorými sa pripravujú textúry a podklady dokumentácie. Je prílohou Systémovej príručky
(kapitola 4.1). Každý súbor začína komentárom s~autorom, verziou, názvom modulu a jeho účelom;
funkcie majú komentár s~popisom, vstupmi a výstupmi; nezávislé časti sú oddelené celoriadkovým
komentárom zo znakov = alebo -. Knižnice \subor{phaser.min.js} a \subor{rot.min.js} sú prevzaté
bez zmeny a nevypisujú sa.

% =============================================================================
\section{Program}

\vypis{\hra}{index.html}{Stránka aplikácie: stôl s~knihou, šablóny dvojstrán, filmová vrstva, údaje počas hry.}
\vypis{\hra}{js/main.js}{Vstupný bod: konfigurácia hry Phaser, kontrola WebGL, spustenie režiséra.}
\vypis{\hra}{js/config.js}{Konštanty hry a zoznamy kresieb.}
\vypis{\hra}{js/story.js}{Texty príbehu (scenár).}
\vypis{\hra}{js/director.js}{Režisér: scenár celej aplikácie, spoločný let kamery knihy a sveta.}
\vypis{\hra}{js/book.js}{Kniha na stole: strany, otáčanie listov, doska, písanie, potvrdenia, kamera knihy.}
\vypis{\hra}{js/cinema.js}{Filmová vrstva: pásy, titulok kapitoly, titulky, zatmenie.}
\vypis{\hra}{js/hud.js}{Údaje počas hry: lístok s~úlohou, zväzky, čas a kroky.}
\vypis{\hra}{js/sound.js}{Zvuková réžia: hudba, prostredie, efekty.}
\vypis{\hra}{js/maze.js}{Logika bludiska: generovanie, voľné bunky, najkratšia cesta, slepé uličky, zväzky.}
\vypis{\hra}{js/iso.js}{Izometrická projekcia a hĺbka pre maliarov algoritmus.}
\vypis{\hra}{js/layout.js}{Geometria sveta okolo bludiska: pásy, strany, rieka, miesta príbehu.}
\vypis{\hra}{js/village.js}{Krajina usadlosti: plán zón, predmety, voda, most, mólo, fontána.}
\vypis{\hra}{js/life.js}{Správanie živých tvorov: chôdza, zvieratá, motýle, rastliny vo vetre.}
\vypis{\hra}{js/scenes/LoadScene.js}{Scéna načítania textúr a zvukov.}
\vypis{\hra}{js/scenes/WorldScene.js}{Scéna sveta: stavba sveta, postava, kamera, rozhovory, oslava.}
\vypis{\hra}{js/scenes/SkyScene.js}{Scéna počasia nad celou obrazovkou: dážď, vietor.}
\vypis{\hra}{css/hra.css}{Vzhľad stola, knihy, filmovej vrstvy a údajov počas hry.}
\vypis{\hra}{spustit.cmd}{Spustenie lokálneho webového servera vo Windows.}

% =============================================================================
\section{Nástroje}

\vypis{\nastroje}{atlas.js}{Príprava atlasu kresieb sveta (npm run atlas).}
\vypis{\nastroje}{ui-textury.js}{Príprava textúr knihy a rozhrania (npm run ui).}
\vypis{\nastroje}{snimky.js}{Snímky obrazovky a ukážkové video do dokumentácie (npm run snimky).}
\vypis{\nastroje}{hlasy.py}{Hlasy rozprávača a postáv – syntéza reči a časy slov (npm run hlasy).}
\vypis{\nastroje}{balik.py}{Zostavenie balíka hry s~tromi príručkami na odovzdanie.}

\end{document}
